\section{Introduction}

Unified multi-task models for facial analysis have emerged as a compelling alternative to task-specific pipelines, enabling shared representations, reduced deployment cost, and real-time inference. FaceXFormer~\cite{narayan2024facexformer} represents the current state of the art, jointly modeling ten heterogeneous facial tasks within a single transformer-based encoder--decoder architecture.

Despite releasing pretrained weights and inference code, the authors do not provide training code, dataset loaders, or multi-task sampling infrastructure, making faithful reproduction non-trivial. This project performs a principled end-to-end reproduction of FaceXFormer, reconstructing the entire training pipeline from scratch.

We treat reproducibility itself as a scientific contribution. Specifically, we: (1) audit paper--code discrepancies, (2) rebuild the architecture and training pipeline, (3) implement all datasets and losses, (4) train multi-task models at scale, and (5) quantify the gap between reported and reproducible performance. Our analysis exposes how underspecified design decisions affect modern vision systems.

\section{Related Work}

\textbf{Task-specific facial models.}
Traditional approaches address each task independently. STARLoss~\cite{zhou2023starloss} reduces ambiguity in landmark detection, TokenHPE~\cite{zhang2023tokenhpe} models head pose via orientation tokens, and ArcFace~\cite{deng2019arcface} defines the dominant paradigm for large-scale face recognition. While effective, these methods require redundant backbone computation.

\textbf{Multi-task face models.}
Early systems such as HyperFace~\cite{ranjan2017hyperface} and All-in-One~\cite{ranjan2017allinone} combine multiple tasks within CNN frameworks. Transformer-based models including SwinFace~\cite{qin2023swinface}, QFace~\cite{sun2024qface}, and Faceptor~\cite{qin2025faceptor} introduce task tokens but remain limited in task coverage and computational efficiency. FaceXFormer achieves full ten-task unification with a lightweight two-layer decoder.

\textbf{Unified vision models.}
Recent work such as CLIP, SAM, and FaRL~\cite{zheng2022farl} demonstrates that shared representations across tasks improve robustness. FaceXFormer extends this paradigm to multi-task facial analysis without large-scale face pretraining.

\textbf{Reproducibility.}
Prior studies highlight systematic gaps between published results and reproducible outcomes. FaceXFormer exemplifies this issue due to missing training infrastructure. We treat this gap as a primary research object.

\section{Approach}

\subsection{Architecture}

We follow FaceXFormer~\cite{narayan2024facexformer}. The encoder is a Swin-B transformer~\cite{liu2021swin} pretrained on ImageNet, producing multi-scale features at strides $4, 8, 16, 32$. These are projected to a common dimension and fused via an MLP-Fusion module (983K parameters) following SegFormer~\cite{xie2021segformer}.

Let $F$ denote the fused face features and $T = \{T_1, \dots, T_n\}$ the learnable task tokens. The FaceX decoder consists of $N=2$ blocks, each performing:

\begin{itemize}
\item \textbf{Task Self-Attention (TSA):} $T \rightarrow T$
\item \textbf{Task-to-Face Cross-Attention (TFCA):} $T \rightarrow F$
\item \textbf{Face-to-Task Cross-Attention (FTCA):} $F \rightarrow T$
\end{itemize}

The final unified head applies an additional TFCA before routing tokens to task heads.

Token counts follow the paper: 19 (segmentation), 68 (landmarks), 9 (head pose), and 1 for each remaining task.

\subsection{Task Heads and Losses}

Each task uses a dedicated head:

\begin{itemize}
\item Segmentation: Dice + CE
\item Landmark: STARLoss
\item Head pose: Geodesic loss
\item Attribute / visibility: BCE with logits
\item Age: $L_1 +$ CE
\item Gender / race / expression: CE
\item Recognition: ArcFace (PartialFC)
\end{itemize}

The joint objective is:

\begin{align}
\mathcal{L} &= \lambda_{\text{seg}} \mathcal{L}_{\text{seg}} 
+ \lambda_{\text{lnd}} \mathcal{L}_{\text{lnd}} 
+ \lambda_{\text{hpe}} \mathcal{L}_{\text{hpe}} \\
&\quad + \lambda_{\text{attr}} \mathcal{L}_{\text{attr}} 
+ \lambda_{\text{age}} \mathcal{L}_{\text{age}} 
+ \lambda_{\text{g/r}} \mathcal{L}_{\text{g/r}} \\
&\quad + \lambda_{\text{exp}} \mathcal{L}_{\text{exp}} 
+ \lambda_{\text{fr}} \mathcal{L}_{\text{fr}} 
+ \lambda_{\text{vis}} \mathcal{L}_{\text{vis}} .
\end{align}

Since $\lambda_i$ are unspecified, we initialize $\lambda_i = 1.0$ and apply gradient-norm balancing.

\subsection{Datasets}

We use all ten datasets from the paper, including CelebAMask-HQ, 300W, BIWI, CelebA, UTKFace, FairFace, RAF-DB, AffectNet, MS1MV3, and COFW.

Images are resized to $224 \times 224$. Augmentation follows Appendix~F: rotation ($\pm 18^\circ$), scaling ($\pm 10\%$), flipping, blur, grayscale, occlusion, and gamma adjustment.

\subsection{Training}

Training uses PyTorch DDP on 8 $\times$ A100 GPUs with FP16. We use AdamW with learning rate $10^{-4}$ and weight decay $10^{-5}$, decayed at epochs 6 and 10 over 12 epochs.

We implement task-balanced sampling via dataset upsampling, ensuring each batch contains all active tasks.

\subsection{Staged Training}

To stabilize optimization, we adopt staged training:

\begin{itemize}
\item Stage 1: 3 tasks (segmentation, landmarks, head pose)
\item Stage 2: 6 tasks (+ attribute, age, gender)
\item Stage 3: 10 tasks (full model)
\end{itemize}

Each stage initializes from the previous checkpoint.

\section{Results}

\subsection{Baseline Verification}

We evaluate the released checkpoint and confirm reported performance on core tasks.

\subsection{Staged Multi-Task Results}

We observe increasing instability as tasks scale. Key reproduced metrics:

\begin{itemize}
\item Segmentation F1: 92.01 $\rightarrow$ 85.91
\item Landmark NME: 4.67 $\rightarrow$ 0.01 (bug)
\item Head pose MAE: 3.52 $\rightarrow$ 0.38 (bug)
\item Attribute accuracy: $\approx$ reproduced
\item Age MAE: 4.17 $\rightarrow$ 35.27 (scale issue)
\item Visibility: 72.56 $\rightarrow$ 99.25 (protocol mismatch)
\end{itemize}

\subsection{Ablations}

We confirm:
\begin{itemize}
\item Bi-directional attention improves stability
\item Balanced sampling is critical
\item Equal $\lambda_i$ leads to loss imbalance
\end{itemize}

\section{Conclusion}

We provide the first full reproduction of FaceXFormer, demonstrating that while the architecture is sound, reproducibility depends critically on undocumented design choices.

Our results highlight three key insights:
\begin{itemize}
\item Loss scale imbalance severely affects multi-task training
\item Dataset and metric inconsistencies dominate performance gaps
\item Missing implementation details prevent faithful reproduction
\end{itemize}

We argue that reproducibility should be treated as a first-class research contribution in modern vision systems.
